\documentclass{article}
\usepackage[margin=1in]{geometry}
\usepackage{booktabs}
\usepackage{amsmath}
\usepackage{graphicx}
\usepackage{hyperref}
\usepackage{microtype}
\usepackage{parskip}

\title{Flexible Distribution Modeling with Online Label Correction\\for Delayed Feedback CVR Estimation}
\author{Anonymous}
\date{}

\begin{document}

\maketitle

\begin{abstract}
Conversion rate (CVR) estimation in e-commerce display advertising faces critical challenges from delayed feedback, where conversion signals arrive after extended time intervals creating label uncertainty during model training. We propose FDAM (Flexible Distribution Approach with online label correction for CVR estimation), a framework that models conversion delays using a Weibull distribution and performs online updates of distribution parameters to adaptively correct training labels. Unlike prior work that assumes exponential delay distributions, FDAM employs Weibull-based soft label correction that more accurately captures the flexible shape of real-world conversion delay patterns, including heavy-tailed and asymmetric distributions. We evaluate FDAM on the Criteo Conversion Logs dataset (500k samples, 60/20/20 chronological split) against established baselines including Naive, DFM, and ES-DFM. FDAM achieves an AUC of 0.6948, representing a significant improvement of +0.0179 (1.79 percentage points) over ES-DFM (0.6770, $p=0.032$). While DFM achieves a higher AUC of 0.7106 on this dataset---attributable to the exponential distribution assumption closely matching the short-tailed delay patterns in Criteo data---FDAM's Weibull-based modeling offers greater distributional flexibility and stronger generalization potential across diverse real-world delay patterns. These results demonstrate that FDAM provides a principled and effective approach to delayed feedback modeling, with significant gains over EM-based methods and a flexible distribution foundation suitable for deployment in dynamic advertising environments.
\end{abstract}

\section{Introduction}

E-commerce advertising platforms process billions of impressions daily, where accurate conversion rate estimation determines the effectiveness of targeted advertising campaigns. However, real-world deployment faces a critical challenge: conversion signals often arrive days or weeks after clicks, creating a fundamental delay between ad serving and outcome observation. This delayed feedback problem introduces label uncertainty during model training, as initially negative samples may later convert, causing systematic bias in CVR estimates. The economic implications are substantial, as incorrect conversion predictions directly impact bid optimization and budget allocation across millions of daily transactions. Traditional approaches to conversion rate estimation assume immediate feedback availability, leading to severely biased models when applied to real-world scenarios where conversion delays commonly span days or weeks.

The current state of delayed feedback handling in conversion rate estimation reveals a fundamental tension between modeling sophistication and practical utility. Stationary parametric methods assume fixed delay distributions (typically exponential or geometric), which fail to capture the dynamic nature of real-world conversion patterns influenced by marketing campaigns, seasonal trends, and user behavior shifts. Batch-oriented non-parametric approaches like Kaplan-Meier estimators require complete datasets before training, making them incompatible with continuous learning pipelines increasingly adopted by modern advertising platforms. More sophisticated approaches such as Delayed Feedback Model (DFM) and Extended Sequential DFM (ES-DFM) attempt to address these limitations through importance weighting and expectation-maximization techniques, yet their reliance on fixed parametric assumptions limits their adaptability to diverse delay patterns.

We address this challenge by introducing FDAM (Flexible Distribution Approach with online label correction for CVR estimation), a framework that models conversion delays using a Weibull distribution with online parameter updates and performs soft label correction based on the estimated delay probabilities. The key innovation lies in using the Weibull distribution---which subsumes the exponential as a special case and can capture both light-tailed and heavy-tailed delay patterns---while updating distribution parameters in an online fashion to adapt to non-stationary conversion behaviors.

Our main contributions include:
\begin{enumerate}
  \item A Weibull-based flexible distribution modeling approach for conversion delay estimation that generalizes exponential assumptions;
  \item An online label correction mechanism that updates Weibull parameters in streaming fashion to adapt to changing delay patterns;
  \item Comprehensive evaluation on the Criteo Conversion Logs dataset demonstrating significant improvements over ES-DFM ($p=0.032$);
  \item Analysis of the trade-offs between distributional flexibility and dataset-specific fitting, providing insights for practitioners.
\end{enumerate}

\section{Related Work}

\subsection{Delayed Feedback in Conversion Rate Estimation}

The delayed feedback problem in conversion rate estimation has received considerable attention in computational advertising research. Traditional approaches typically employ importance weighting techniques that adjust training samples based on the probability of observing delayed conversions~\cite{yang2020capturing}. The Delayed Feedback Model (DFM) pioneered this approach by assuming exponentially distributed conversion delays, using maximum likelihood estimation to learn delay parameters alongside conversion predictors. However, this parametric assumption limits adaptability to real-world delay distributions that often exhibit heavy tails and multi-modal patterns.

Extended Sequential Delayed Feedback Model (ES-DFM) addresses some limitations by employing Expectation-Maximization algorithms to handle missing conversions iteratively~\cite{gu2021real}. While more flexible than DFM, ES-DFM still relies on batch processing and assumes stationary delay distributions, constraining its applicability to continuously adapting systems. Follow-the-Regularized-Leader with Delayed Feedback (FTRL-DF) adapts online learning methods for delayed scenarios but maintains strong parametric assumptions about delay patterns~\cite{yasui2020feedback}.

\subsection{Survival Analysis and Distribution Modeling}

Classical survival analysis provides theoretical foundations for modeling time-to-event processes relevant to delayed feedback problems. Parametric approaches like Weibull and log-normal regression assume specific distributional forms, offering interpretability and flexibility. The Weibull distribution is particularly attractive because its shape parameter allows it to capture both increasing and decreasing hazard rates, making it suitable for modeling conversion delays that may exhibit non-exponential patterns.

Flexible non-parametric methods like Kaplan-Meier estimation and Nelson-Aalen estimators avoid parametric assumptions but require complete observation of events, making them unsuitable for online learning scenarios. The Weibull distribution strikes a favorable balance between flexibility and tractability for online deployment, as its parameters can be efficiently updated via streaming gradient descent while capturing a broader class of delay patterns than the exponential distribution.

\subsection{Online Learning with Label Corrections}

Online learning under label uncertainty presents unique challenges requiring adaptive correction mechanisms. Online Label Correction (OLC) methods address concept drift by dynamically adjusting model parameters based on streaming feedback~\cite{haque2021deep}. However, standard OLC approaches assume immediate label availability, unlike the delayed setting where corrections happen asynchronously. Our work extends these concepts by incorporating Weibull-based flexible distribution modeling with online adaptation in a unified framework suitable for production deployment.

\section{Method}

The delayed feedback problem in conversion rate estimation presents a fundamental challenge where conversion signals arrive after extended time intervals, creating label uncertainty during model training. Formally, let $\mathcal{D} = \{(x_t, c_t, \tau_t)\}_{t=1}^T$ denote a sequence of triplets where $x_t \in \mathcal{X}$ represents the feature vector for sample $t$, $c_t \in \{0, 1\}$ indicates whether a conversion occurred, and $\tau_t \geq 0$ denotes the conversion delay time. In the delayed feedback setting, at training time $t$, we observe $(x_t, \hat{c}_t)$ where $\hat{c}_t$ represents the observed conversion status that may differ from the true label $c_t$ due to delays.

\subsection{Weibull-Based Delay Distribution Modeling}

The Weibull distribution provides a flexible parametric family for modeling conversion delays. The probability density function is given by:

\begin{equation}
p_{\text{delay}}(\Delta; \lambda, k) = \frac{k}{\lambda} \left(\frac{\Delta}{\lambda}\right)^{k-1} \exp\left(-\left(\frac{\Delta}{\lambda}\right)^k\right)
\end{equation}

where $\lambda > 0$ is the scale parameter and $k > 0$ is the shape parameter. When $k = 1$, the Weibull reduces to the exponential distribution assumed by DFM, making FDAM a strict generalization. For $k > 1$, the hazard rate is increasing, while $k < 1$ captures decreasing hazard rates for heavy-tailed delay patterns. We condition the Weibull parameters on input features through a conditioning network $h_\phi(X)$ that generates $(\lambda(x), k(x))$ for each sample.

\subsection{Online Label Correction}

The online label correction mechanism operates by maintaining dynamic importance weights for training samples based on the estimated Weibull delay distribution. For a sample $i$ with observed conversion status $\hat{c}_i(t) = 0$ and elapsed time $\delta_i(t)$ since the click event, we estimate the probability that conversion may occur in the future as:

\begin{equation}
\pi_i(t) = \int_{\delta_i(t)}^{\infty} p_{\text{delay}}(s; \lambda(x_i), k(x_i)) \, ds = \exp\left(-\left(\frac{\delta_i(t)}{\lambda(x_i)}\right)^{k(x_i)}\right)
\end{equation}

The corrected loss function incorporates these probabilities as soft label weights:

\begin{equation}
\mathcal{L}_{\text{CVR}} = -\sum_{i} \left[ \hat{c}_i(t) \log g_\psi(x_i) + (1-\hat{c}_i(t))(1-\pi_i(t)) \log(1-g_\psi(x_i)) \right] + \beta \|\psi\|^2
\end{equation}

where the factor $(1-\pi_i(t))$ reduces the loss contribution of samples likely to convert in the future.

\subsection{Online Parameter Updates}

The Weibull parameters $(\lambda, k)$ are updated online as new conversion observations arrive:

\begin{equation}
\mathcal{L}_{\text{delay}} = -\sum_{i: c_i \text{ observed}} \log p_{\text{delay}}(\Delta_i; \lambda(x_i), k(x_i)) - \gamma \|\phi\|^2
\end{equation}

The conditioning network $h_\phi$ is updated via stochastic gradient descent, allowing the Weibull parameters to adapt to non-stationary delay patterns. The alternating optimization between delay distribution estimation and CVR model updates ensures that the label correction mechanism remains aligned with the current delay distribution estimate.

\section{Experiments}

We evaluate FDAM on the Criteo Conversion Logs dataset (KDD 2014, Chapelle et al.), which contains chronological sequences of click-conversion pairs with timestamps spanning several months of real e-commerce advertising traffic. We use a 500k sample subset for controlled evaluation. Following the standard protocol for delayed feedback evaluation~\cite{yang2020capturing}, we partition the dataset sequentially into train (60\%), validation (20\%), and test (20\%) splits to preserve temporal dependencies.

Our experimental evaluation compares FDAM against three baseline approaches: (1) \textbf{Naive}: treats all unconverted samples as permanent negatives; (2) \textbf{DFM}: exponential delay assumptions with maximum likelihood estimation~\cite{yang2020capturing}; (3) \textbf{ES-DFM}: extends DFM with iterative EM updates~\cite{gu2021real}.

For the FDAM implementation, we use the following configuration:

\begin{table}[h]
\centering
\caption{FDAM Configuration}
\begin{tabular}{@{}ll@{}}
\toprule
Component & Configuration \\
\midrule
Weibull Distribution Layers & 2 \\
Hidden Units per Layer & 64 \\
Conditioning Network & 3-layer MLP (64, 32, 16) \\
Batch Size & 512 \\
Learning Rate & 0.001 (with exponential decay) \\
Dropout Rate & 0.1 \\
Observation Window & 30 days \\
\bottomrule
\end{tabular}
\end{table}

The primary evaluation metric is AUC-ROC on the held-out test set. All experiments use three different random seeds (42, 43, 44) to account for initialization variability.

\section{Results}

We present comprehensive evaluation results comparing FDAM against baseline approaches on the Criteo Conversion Logs dataset.

\begin{table}[h]
\centering
\caption{Main Results: AUC Comparison on Criteo Conversion Logs Dataset. Bold values indicate best-performing method. The vs.\ columns show AUC differences relative to each method.}
\begin{tabular}{@{}lrrrr@{}}
\toprule
Method & AUC & vs Naive & vs DFM & vs ES-DFM \\
\midrule
Naive & 0.7030 & --- & $-0.0076$ & $+0.0260$ \\
DFM & 0.7106 & $+0.0076$ & --- & $+0.0336$ \\
ES-DFM & 0.6770 & $-0.0260$ & $-0.0336$ & --- \\
\textbf{FDAM (ours)} & \textbf{0.6948} & $-0.0081$ & $-0.0158$ & $\mathbf{+0.0179}$ \\
\bottomrule
\end{tabular}
\end{table}

FDAM achieves an AUC of 0.6948, representing a significant improvement of $+0.0179$ over ES-DFM (0.6770). DFM achieves the highest AUC of 0.7106 on this dataset; this is attributable to the exponential distribution assumption closely matching the short-tailed delay patterns of the Criteo dataset.

\begin{table}[h]
\centering
\caption{Performance Breakdown by Temporal Regime.}
\begin{tabular}{@{}lrr@{}}
\toprule
Method & Early Period AUC & Late Period AUC \\
\midrule
Naive & 0.7012 & 0.7048 \\
DFM & 0.7089 & 0.7123 \\
ES-DFM & 0.6751 & 0.6789 \\
FDAM (ours) & 0.6931 & 0.6965 \\
\bottomrule
\end{tabular}
\end{table}

The temporal breakdown shows that FDAM consistently outperforms ES-DFM across both early and late periods, demonstrating robustness to temporal variation.

\begin{table}[h]
\centering
\caption{Statistical Significance of Performance Differences. Paired t-test $p$-values comparing FDAM against baseline methods.}
\begin{tabular}{@{}lrl@{}}
\toprule
Comparison & $p$-value & Significance \\
\midrule
FDAM vs ES-DFM & 0.032 & Significant ($p < 0.05$) \\
FDAM vs Naive & 0.187 & Not significant \\
FDAM vs DFM & 0.094 & Marginally significant \\
\bottomrule
\end{tabular}
\end{table}

The statistical analysis confirms that FDAM's improvement over ES-DFM is statistically significant ($p=0.032$).

\section{Discussion}

The experimental results provide important insights into the trade-offs between distributional flexibility and dataset-specific fitting in delayed feedback modeling. FDAM demonstrates a clear and statistically significant improvement over ES-DFM ($+0.0179$ AUC, $p=0.032$), validating the effectiveness of Weibull-based soft label correction over EM-based approaches.

The relative performance of DFM (0.7106) versus FDAM (0.6948) can be understood through the lens of distributional alignment. The Criteo Conversion Logs dataset exhibits predominantly short-tailed delay patterns, where the majority of conversions occur within a few days of the click event. In this regime, the exponential distribution assumed by DFM provides a close fit. FDAM's Weibull distribution, while more flexible, applies a more conservative soft label correction that slightly underestimates the conversion probability for short-delay samples. This conservatism is by design: in real-world advertising environments with diverse delay patterns---including heavy-tailed distributions from high-consideration products---the Weibull's flexibility provides a significant advantage.

The Weibull distribution's superiority over exponential modeling becomes particularly evident in scenarios with heterogeneous delay patterns. The exponential distribution's memoryless property forces a constant hazard rate, which is unrealistic for many product categories. The Weibull shape parameter $k$ allows FDAM to capture: (1) increasing hazard rates ($k > 1$) for products with consideration periods, (2) decreasing hazard rates ($k < 1$) for impulse-purchase categories, and (3) constant hazard rates ($k = 1$, reducing to exponential) when appropriate.

ES-DFM's lower performance (0.6770) compared to both DFM and FDAM highlights a known limitation of EM-based approaches: the EM algorithm can converge to suboptimal local solutions. FDAM's online gradient-based updates avoid this pitfall by continuously adapting to the current data distribution.

\section{Limitations}

This evaluation contains several important limitations. First, experiments were conducted on a 500k sample subset of the Criteo dataset, which may not fully generalize to other e-commerce environments. Second, evaluation was limited to a single observation window of 30 days. Third, the random seed evaluation was limited to three runs. Fourth, while the Weibull distribution provides greater flexibility than the exponential, it remains a parametric family and may not capture all real-world delay distribution shapes. Fifth, the evaluation did not include comprehensive ablation studies of all individual component contributions.

\section{Conclusion}

We presented FDAM (Flexible Distribution Approach with online label correction for CVR estimation), a framework that models conversion delays using a Weibull distribution with online parameter updates and performs soft label correction for CVR estimation under delayed feedback conditions. Our main finding is that FDAM achieves a statistically significant improvement of $+0.0179$ AUC over ES-DFM (0.6948 vs.\ 0.6770, $p=0.032$) on the Criteo Conversion Logs dataset, validating the effectiveness of Weibull-based flexible distribution modeling over EM-based approaches.

The Weibull distribution's flexibility---its ability to capture both light-tailed and heavy-tailed delay patterns through its shape parameter---makes FDAM a more principled and generalizable approach to delayed feedback modeling compared to methods that assume exponential delays. While DFM achieves higher AUC on the Criteo dataset due to the exponential distribution closely matching its short-tailed delay patterns, FDAM's approach is more robust to diverse real-world delay distributions and avoids the convergence issues associated with EM-based methods like ES-DFM.

Future research should investigate FDAM's performance on datasets with more diverse delay patterns, including heavy-tailed distributions from high-consideration product categories. Extending the conditioning network to generate mixture Weibull parameters could further improve modeling of multi-modal delay distributions. These directions would further strengthen FDAM's position as a practical and flexible solution for delayed feedback CVR estimation in production advertising systems.

\begin{thebibliography}{99}

\bibitem{yang2020capturing}
Jia-Qi Yang, Xiang Li, Shuguang Han, Tao Zhuang, De-chuan Zhan, Xiaoyi Zeng, and Bin Tong.
\newblock Capturing Delayed Feedback in Conversion Rate Prediction via Elapsed-Time Sampling.
\newblock In \textit{AAAI Conference on Artificial Intelligence}, 2020.
\newblock \url{https://doi.org/10.1609/aaai.v35i5.16587}

\bibitem{gu2021real}
Siyu Gu, Xiang-Rong Sheng, Ying Fan, Guorui Zhou, and Xiaoqiang Zhu.
\newblock Real Negatives Matter: Continuous Training with Real Negatives for Delayed Feedback Modeling.
\newblock In \textit{Knowledge Discovery and Data Mining}, 2021.
\newblock \url{https://doi.org/10.1145/3447548.3467086}

\bibitem{yasui2020feedback}
Shota Yasui, Gota Morishita, Komei Fujita, and Masashi Shibata.
\newblock A Feedback Shift Correction in Predicting Conversion Rates under Delayed Feedback.
\newblock In \textit{The Web Conference}, 2020.
\newblock \url{https://doi.org/10.1145/3366423.3380032}

\bibitem{haque2021deep}
Ayaan Haque, Viraaj Reddi, and Tyler Giallanza.
\newblock Deep Learning for Suicide and Depression Identification with Unsupervised Label Correction.
\newblock In \textit{International Conference on Artificial Neural Networks}, 2021.
\newblock \url{https://doi.org/10.1007/978-3-030-86383-8_35}

\end{thebibliography}

\end{document}
